\scanpage{57}
\mspage{28}
se grouperaient donc \inserted{classes de} \deleted{\unsure{par}\ \unsure{des}} paquets
(en regardant les centralisateurs des ss-groupes
images de $\Gamma^\circ$ par ses sections,\deleted{)} et en
mettant ensemble des scindages qui ont des
$\unsure{\mathrm{Zentr}|\mathrm{Centr}}(\Gamma^\circ)$ conjugués (même s'ils ne sont
eux-mêmes conjugués). Donc on retrouve
encore une description de $X(K)$ (donc
aussi des $X(K')$ pour tt ext. finie \inserted{$K'$} de $K$)
en termes de l'extension $E_L$ de $E_K$ par
$\unsure{\Pi|\pi}_{L/K}$, en $=$ temps qu'une façon de
reconstituer les $U=X\setminus\bar I$\ldots
Donc en fait c'est la structure $E_L\unsure{-|\longrightarrow}E_K$
qui est la plus riche a priori, et de
loin plus commode pour le genre $0$ et $1$,
où la considération des $U$ de type $(g,\nu)$\unsure{.|,}
$(2g+\nu\geq3)$ casse le groupe ``continu''
d'automorphismes\ldots La forme ``modélique''
de la conjecture précédente revient
à la forme ``birationnelle'', quand on y
précise cette dernière en disant que tt
\margin{56}

